\documentclass[11pt,letterpaper]{article}
\usepackage[top=0.7in,bottom=0.7in,left=0.75in,right=0.75in]{geometry}
\usepackage[hidelinks]{hyperref}
\usepackage[T1]{fontenc}
\usepackage[utf8]{inputenc}
\renewcommand{\familydefault}{\sfdefault}
\pagestyle{empty}
\setlength{\parindent}{0pt}
\setlength{\parskip}{12pt}
\linespread{1.12}
\begin{document}
\raggedright
{\Large\textbf{Candidate}}\\[4pt]
San Francisco, CA \enspace|\enspace \href{mailto:candidate@example.com}{candidate@example.com} \enspace|\enspace \href{https://portfolio.example.com}{portfolio.example.com}

\vspace{5pt}
\textbf{Brain Co. | AI Platform Engineer Co-Op}

I'm interested in Brain Co. because you're applying AI to practical problems like construction permitting and hospital workflows. I want to help build the infrastructure that lets your engineers deploy changes confidently and keep those applications running well. Seeing that work support real customer needs is a big part of the appeal for me.

Your focus on developer tooling, automation, and AI workload performance fits the work I've enjoyed most. Building an internal AI platform at Paychex taught me to think about what other engineers need to use a system effectively. Working on production AI agents at Real Wave made latency and operating costs practical engineering concerns for me. I'd bring that perspective to your platform team while developing deeper experience with Kubernetes, Terraform, and observability.

I also like that this co-op combines ownership with mentorship through code reviews and design discussions. I'm a computer science student at RIT and am available for both Spring and Summer 2027, including a role spanning both terms. I'd welcome the chance to contribute across that longer period.
\end{document}
